\scanpage{175}
\mspage{83}
on a loc. une décomposition en produits de fibrés en droites, d'où (en prenant le produit des fibrés principaux associés) un ouvert
\[
N^*\subset N
\]
\[
(289)\qquad\left\{\begin{array}{ll}
\text{en particulier}&N^*_{d_*}\subset N_{d_*}\\
&N^*_{d_*,d_i}\subset N_{d_*,d_i}\\
&N^*_{d_*;(d'_0,d''_0)}\subset N_{d_*;(d'_0,d''_0)}
\end{array}\right.
\]
qui est d'ailleurs un torseur sous un \deleted{fibré} tore tordu convenable, que je me dispense d'expliciter.

Considérons alors
\[
(290)\qquad\check N^{d'_0}_{d_*,d''_0}=\text{image inverse de }\check N^*_{d_*;(d'_0,d''_0)}
\]
dans $\check N_{d_*,d''_0}$, par l'hom
\diagram{\begin{tikzcd}[ampersand replacement=\&]
\check N_{d_*,d''_0}\arrow[r,no head]\&\check N_{d_*;(d'_0,d''_0)}
\end{tikzcd}}
de (287)

de sorte qu'on a
\[
(291)\qquad\check N^{d'_0}_{d_*,d''_0}\quad\text{est fibré sur }\overset{\substack{\check N^*\\\Vert}}{\check N^*_{d_*;(d'_0,d''_0)}},
\]
et est un torseur sous le fibré vectoriel sur $\check N^*$, image inverse de $\check N_{d_*,\unsure{d_0|d'_0}}$ par $\check N^*\to D_{d^{\unsure{?}}_*}$.
\margin{174}
